\subsection{Data Extraction and Visualization Architecture}

To capture the dynamic evolution of the socio-ecological system and mathematically evaluate the efficacy of the tested policies, the simulation employs a dual-tiered data collection architecture. This separates real-time spatial validation from rigorous, publication-quality quantitative telemetry.

\paragraph{Part 1: Real-Time Spatial Visualization (The Interactive Server)}
For qualitative observation and stakeholder demonstration, the model utilizes the \texttt{mesa.visualization} module to host a local \texttt{ModularServer}. The graphical interface features a custom Portrayal Factory that partitions the environment into seven distinct \texttt{CanvasGrid} instances, allowing for real-time spatial tracking of Household and Enforcement agents across the different barangays. 

To monitor macro-level compliance concurrently, the server utilizes a \texttt{ChartModule} locked to a high-resolution 1000px canvas. While this User Interface (UI) is highly effective for observing localized spatial clustering and the immediate visual impact of the Mayor Agent's quarterly interventions, browser-based data collectors are constrained by caching limits and UI-thread processing delays, making them insufficient for final statistical extraction.

\begin{figure}[h]
    \centering
    % Ensure you have uploaded the file 'simulation_ui.jpg' to your project
    \includegraphics[width=0.8\linewidth]{simulation_ui.jpg}
    \caption{The Bacolod Multi-View Simulation UI (Interactive Validation Environment)}
    \label{fig:simulation_ui}
\end{figure}

\paragraph{Part 2: Headless Execution and Deterministic Telemetry}
To generate the empirical results for policy comparison, the study strictly bypasses the web UI and executes deterministic, headless simulation loops via automated benchmarking scripts (e.g., \texttt{compare\_algo.py}). 

Running the model in headless mode eliminates browser caching inaccuracies. During these benchmarking runs, the script actively queries the memory states of all instantiated \texttt{HouseholdAgent} objects at every temporal tick, calculating the exact floating-point global compliance rate. To ensure publication-quality data visualization, the script utilizes the \texttt{matplotlib.use('Agg')} backend renderer. Rather than relying on screen captures, the system algorithmically generates and serializes exact vector graphics (exported at 300 DPI resolution) mapping the temporal trajectories of the ``Status Quo,'' ``Pure Incentives,'' ``Pure Enforcement,'' and the ``HuDRL Mayor Agent'' policy regimes.

\paragraph{Part 3: DRL Training Logs and Sensitivity Outputs}
Beyond standard ABM compliance tracking, the architecture integrates specialized telemetry pipelines for the artificial intelligence and parameter validation components:
\begin{itemize}
    \item \textbf{TensorBoard Integration:} During the parallelized training phase of the Deep Reinforcement Learning (DRL) agent, the \texttt{SubprocVecEnv} is wrapped in a \texttt{VecMonitor}. This continuously streams the agent's algorithmic metrics—specifically the episodic reward convergence, policy entropy, and value loss—directly to local TensorBoard logs, providing a transparent record of the neural network's learning curve \citep{Raffin2021}.
    \item \textbf{Sobol Statistical Extraction:} For the variance-based Global Sensitivity Analysis (detailed in Section 3.3.1), the \texttt{sobol\_sensitivity.py} script executes hundreds of parameterized environments in the background. It utilizes the \texttt{SALib} library to compute the First-Order ($S_1$) and Total-Order ($S_T$) sensitivity arrays, outputting the exact statistical indices to the console and rendering a dedicated comparative bar chart.
\end{itemize}

This automated, multi-tiered extraction pipeline mathematically guarantees that all comparative analyses, sensitivity indices, and learning curves presented in the subsequent results are derived from exact computational states rather than graphical approximations \citep{McCulloch2022}.